% =========================
% report/system_overview.tex
% =========================

\subsection{خلاصهٔ خط لوله}
سامانه با یک خط لولهٔ استاندارد ثبت تصویر مبتنی بر ویژگی کار می‌کند:
\begin{enumerate}
  \item \textbf{پیش‌پردازش} $\mathcal{I}_T$ و $\mathcal{I}_Q$ (خاکستری‌سازی، کنترل اندازه، و در صورت نیاز افزایش کنتراست).
  \item \textbf{استخراج ویژگی} (نقاط کلیدی و توصیف‌گر) با ORB \cite{Rublee2011ORB}.
  \item \textbf{تطبیق توصیف‌گرها} با فاصلهٔ همینگ و پالایش تطبیق‌ها (Cross-check و/یا فیلتر فاصله).
  \item \textbf{تخمین هموگرافی} $\mathbf{H}$ به کمک RANSAC برای حذف پرت‌ها \cite{FischlerBolles1981}.
  \item \textbf{وارپ} کردن الگو به چارچوب عکس با $\mathbf{H}$ و تولید خروجی هم‌پوشانی.
  \item \textbf{خروجی‌گیری} از معیارهای کمی و تصاویر دیباگ برای ارزیابی علمی.
\end{enumerate}

\subsection{ساختار پوشه‌ها (پیشنهادی)}
گزارش فرض می‌کند پروژه ساختار زیر را دارد:
\begin{verbatim}
report/
  main.tex
  *.tex
  references.bib
src/
  (ماژول‌های پایتون)
data/
  templates/
  photos/
  annotations/   (اختیاری)
results/
  overlays/
  figures/
  metrics/
  matches/       (اختیاری)
\end{verbatim}

\subsection{چرایی انتخاب روش مبتنی بر ویژگی}
این رویکرد \textbf{بدون آموزش} و با هزینهٔ محاسباتی پایین اجرا می‌شود و اگر تصویر حاوی ساختار/بافت کافی باشد، ORB+RANSAC می‌تواند هم‌ترازی قابل قبولی ارائه کند \cite{Rublee2011ORB,FischlerBolles1981}.
